\documentclass[10pt,a4paper]{amsart}
\usepackage[margin=19mm]{geometry}
\usepackage{amsmath,amssymb,mathtools}
\usepackage[scr=rsfs,cal=euler]{mathalfa}
\usepackage{xfrac}
\usepackage[unicode,hidelinks,bookmarks=false]{hyperref}
\newcommand{\ie}{i.e.\ }
\newcommand{\cf}{cf.\ }
\newcommand{\resp}{respectively\ }
\newcommand{\Subsection}{\subsection}
\providecommand{\nobreakdash}{\nobreak\hbox{-}}
\setlength{\emergencystretch}{2em}
\multlinegap0pt
\newcommand{\R}{\mathbb R}
\theoremstyle{plain}
\newtheorem{theorem}{Theorem}[section]
\newtheorem{lemma}[theorem]{Lemma}
\newtheorem{proposition}[theorem]{Proposition}
\newtheorem{corollary}[theorem]{Corollary}
\theoremstyle{remark}
\newtheorem{remark}[theorem]{Remark}
\title[A Smooth Nonnegatively Curved Sphere Whose Zero-Curvature Point Locus Has Hausdorff Dimension 1/2]{A Smooth Nonnegatively Curved Sphere Whose Zero-Curvature Point Locus Has Hausdorff Dimension 1/2}
\author{Yuhang Liu}
\date{Source: arXiv:2609.27898v1 (CC BY 4.0)}
\begin{document}
\maketitle

\noindent\emph{Source and licence.} The delivered work is the article shown in the title above, version arXiv:2609.27898v1 (CC BY 4.0), distributed under the Creative Commons Attribution 4.0 International licence, as recorded in the arXiv licence metadata for this exact version and shown on the abstract page saved with this item. The full licence notice, which carries the licence URL and the address of that abstract page, is the bundled file \texttt{licenses/arxiv-2609.27898-CC-BY-4.0.txt}.
\par\smallskip

\noindent\textit{Editorial scope.} Counted unit: the article's Proposition (source label prop:local-model) (source0.tex lines 217--223, source label \texttt{prop:local-model}): ``The function is convex on . Its Hessian is positive definite on , and at every point of its Hessian has rank exactly . Moreover, 0 U 0(x) 3 3 100 (x B m).''. It is delivered together with the article's complete original proof (source0.tex lines 225--257, one proof environment adjacent to the statement, 33 lines), copied line for line from the article's own text. Required context retained uncounted: eq:u-properties (lines 201--205). Editorial formatting only: an AMS \texttt{amsart} preamble that restates the article's own macro definitions and its own theorem declarations, and the theorem environments the retained lines use; the article's own class file is neither shipped nor input; the retained environments are renumbered sequentially in order of appearance, whereas the article prints them with its own numbering; the article's equation numbering is not reproduced and every cross-reference inside the retained text resolves to the wrapper's numbering. No citation occurs inside any retained line, so the wrapper carries no bibliography; All retained bodies are exact copies of the bundled \texttt{sources/211523/source0.tex}, so no omission receipt applies. No asset-inclusion command occurs inside any retained span and the package ships no image asset.


\section*{Retained context: eq:u-properties (source0.tex lines 201--205)}

\begin{equation}\label{eq:u-properties}
 u''(t)=\varepsilon q(t),\qquad u(0)=u'(0)=0,
 \qquad 0\leq u(t)\leq\frac\varepsilon2t^2
 \quad (|t|\leq1).
\end{equation}


\section{The counted unit: Proposition (source label prop:local-model) and its complete original proof}

\begin{proposition}\label{prop:local-model}
The function $U$ is convex on $B^m$.  Its Hessian is positive definite on $B^m\setminus K$, and at every point of $K$ its Hessian has rank exactly $m-1$.  Moreover,
\begin{equation}\label{eq:U0-bound}
 0\leq U_0(x)\leq3\varepsilon=\frac3{100}
 \qquad (x\in B^m).
\end{equation}
\end{proposition}

\begin{proof}
Since the constant $-1/2$ does not affect the Hessian, it is enough to consider $U_0$.  In the splitting $\R\oplus\R^{m-1}$,
\begin{equation}\label{eq:hessian-U}
 D^2U_0(t,y)
 =\varepsilon
 \begin{pmatrix}
 q(t)+|y|^2 & 2t y^{T}\\
 2t y & (4+t^2)I_{m-1}
 \end{pmatrix}.
\end{equation}
The lower-right block is positive definite.  Its Schur complement is
\begin{align}
 \varepsilon\left(q(t)+|y|^2
 -\frac{4t^2}{4+t^2}|y|^2\right)
 &=\varepsilon\left(
 q(t)+\frac{4-3t^2}{4+t^2}|y|^2\right).\label{eq:schur}
\end{align}
Because $|t|<1$, the coefficient $(4-3t^2)/(4+t^2)$ is positive.  Thus the Schur complement is nonnegative, and it vanishes exactly when $q(t)=0$ and $y=0$, namely at the points of $K$.  This proves convexity and positive definiteness off $K$.  At a point $(t,0)\in K$, the matrix in \eqref{eq:hessian-U} is
\[
 \varepsilon
 \begin{pmatrix}
 0&0\\0&(4+t^2)I_{m-1}
 \end{pmatrix},
\]
which has rank $m-1$.

Finally, \eqref{eq:u-properties} and $|t|,|y|<1$ give
\[
 0\leq U_0(t,y)
 \leq\frac\varepsilon2+\frac{5\varepsilon}{2}
 =3\varepsilon.
\]
\end{proof}

\end{document}
